\begin{definition}
\label{definition-quasi-finite}
\begin{reference}
\cite[II Definition 6.2.3]{EGA}
\end{reference}
Let $f : X \to S$ be a morphism of schemes.
\begin{enumerate}
\item We say that $f$ is {\it quasi-finite at a point $x \in X$}
if there exist an affine neighbourhood $\Spec(A) = U \subset X$
of $x$ and an affine open $\Spec(R) = V \subset S$ such that
$f(U) \subset V$, the ring map $R \to A$ is of finite type,
and $R \to A$ is quasi-finite at the prime of $A$ corresponding to $x$
(see above).
\item We say $f$ is {\it locally quasi-finite} if $f$ is
quasi-finite at every point $x$ of $X$.
\item We say that $f$ is {\it quasi-finite} if $f$ is of finite type
and every point $x$ is an isolated point of its fibre.
\end{enumerate}
\end{definition}